% 一系列主题logo

\newcommand{\liuyin}[1][5ex]{%
	\adjustbox{valign=m, raise=#1, height=0.8em}{\input{ysy-latex/theme/liuyin_icon.pdf_tex}}%
}

%=============================================================================

% 以下设计来自ElegantBook

% tikz绘制UI
\usetikzlibrary{backgrounds,calc,shadows,positioning,fit}

% 定义带圈数字命令
\newcommand*\circled[1]{\tikz[baseline=(char.base)]{
		\node[shape=circle,draw,inner sep=1pt] (char) {#1};}}
% 作用：生成带圆圈的数字（如⑪），常用于强调序号。
% 参数：#1 为要包围的内容（如数字）。
% 关键属性：
% inner sep=1pt：控制圆圈与内容的间距。
% baseline：对齐文本基线，确保与文字垂直居中。
% 示例：\circled{5} 生成带圈数字⑤。

% 定义三种不同层级的无序列表符号：
\newcommand*{\eitemi}{\tikz \draw [baseline, ball color=MainColor,draw=none] circle (3pt);}
\newcommand*{\eitemii}{\tikz \draw [baseline, fill=MainColor,draw=none,circular drop shadow] circle (3pt);}
\newcommand*{\eitemiii}{\tikz \draw [baseline, fill=MainColor,draw=none] circle (3pt);}
% 效果：
% \eitemi：立体球状符号（ball color）。
% \eitemii：带阴影的实心圆点。
% \eitemiii：普通实心圆点。

%=============================================================================

% 用户信息

% -- 用户可设置封面项 --
\newcommand{\TitleName}{}
\newcommand{\AuthorName}{}
\newcommand{\Address}{}
\newcommand{\Preface}{}
\newcommand{\CoverIllustration}{}
\newcommand{\AddInfo}{}
\newcommand{\Abstract}{}